\chapter{Kademeli Çözüm ve Değerlendirme Rehberi}
\label{app:answer-key}

Bu ek, kitabın temel konularına ilişkin seçilmiş alıştırmalar için kademeli
çözüm desteği sunar. Konu kümeleri, iki eski kitap projesinden korunan ve
birleşik yapıya aktarılan soru havuzunu izlediği için her bölümdeki bütün
etkinliklerle birebir bir cevap listesi değildir. Her maddede önce çözümü
açmadan ilerlemeyi sağlayan bir \emph{ipucu}, ardından kısa \emph{yanıt veya
değerlendirme ölçütü} verilir. Açık uçlu sorularda örnekler tek doğru
değildir; belirtilen temel kavramları doğru ilişkilendiren başka yanıtlar da
kabul edilebilir. Ayrıntılı hesaplarda ara yuvarlama yerine son aşamada
yuvarlama yapılmalıdır.

\section{İstatistiksel düşünme ve araştırma süreci}

\begin{enumerate}[leftmargin=*,itemsep=5pt]
\item \emph{İpucu:} Soruyu kim, ne ve hangi sonuç biçiminde ayırın.
\textbf{Ölçüt:} Hedef evren, tek gözlem birimi, açıklayıcı değişken ve yanıt
değişkeni birbiriyle tutarlı biçimde tanımlanmalıdır.
\item \emph{İpucu:} Sayı biçimine değil, değerlerin taşıdığı anlama bakın.
\textbf{Yanıt:} Okul türü nominal; yaş ve sınav süresi oran ölçeğinde sürekli;
öğrenci numarası nominal tanımlayıcı; başarı sırası sıralıdır. Sırasıyla
frekans, ortalama/medyan, frekans, medyan/yüzdelik ve ortalama/medyan uygundur.
\item \emph{İpucu:} Evren özelliği ile örneklem hesabını ayırın.
\textbf{Yanıt:} İlçedeki bütün öğrencilerin ortalama okuma süresi $\mu$
parametresi; seçilen öğrencilerin ortalaması $\bar x$ istatistiğidir.
\item \emph{İpucu:} Ankete kimlerin eriştiğini ve kimlerin yanıt vermeyi
seçtiğini düşünün. \textbf{Yanıt:} Kapsam yanlılığı ve gönüllülük/öz-seçim
yanlılığı temel iki risktir.
\item \emph{İpucu:} Küçük fakat önemli alt grupları düşünün.
\textbf{Ölçüt:} Örneğin fakülte veya program düzeylerinin ayrı temsil edilmesi
gerektiğinde her tabakadan olasılıklı seçim önerilmelidir.
\item \emph{İpucu:} Her iki değişkeni birlikte etkileyebilecek nedenleri arayın.
\textbf{Ölçüt:} Önceki başarı, motivasyon, aile desteği, uyku veya çalışma
alışkanlığı gibi en az iki makul karıştırıcı verilmelidir.
\item \emph{İpucu:} Satır sayısı ile bağımsız öğrenci sayısını ayırın.
\textbf{Yanıt:} Aynı öğrenci iki satırla temsil edildiğinden gözlemler bağımsız
değildir; analiz birimi sınav girişi seçilirse öğrenci içi kümelenme modellenmelidir.
\item \emph{İpucu:} Sıfırın gerçek bir ölçüm mü yoksa kod mu olduğunu sorun.
\textbf{Yanıt:} Kod, gerçek sıfır gibi hesaba katılırsa ortalama düşer ve yapay
bir yığılma oluşur; eksikler ayrı eksik değer olarak tanımlanmalıdır.
\item \emph{İpucu:} Rastgele hata ile sistematik kapsam hatasını ayırın.
\textbf{Ölçüt:} Örneğin yalnız çevrim içi platform kullanıcılarından gelen çok
büyük örneklemin, platformu kullanmayanları temsil etmeyeceği açıklanmalıdır.
\item \emph{İpucu:} Süreci toplama--temizleme--analiz--raporlama olarak kurun.
\textbf{Ölçüt:} Hedef evren ve değişkenler, onam/gizlilik, eksik ve aykırı
değer kuralları, grafiksel inceleme, uygun çıkarım ve sınırlılıklar bulunmalıdır.
\end{enumerate}

\section{Evren, örneklem ve veri türleri}

\begin{enumerate}[leftmargin=*,itemsep=5pt]
\item \emph{İpucu:} Sonucun genelleneceği grup ile erişilebilen listeyi ayırın.
\textbf{Ölçüt:} Örneğin bütün lise öğrencileri hedeflenirken yalnız devlet
liselerinin kayıtlarına erişilmesi uygun bir örnektir.
\item \emph{İpucu:} Bilinmeyen evren sayısını ve onun örneklem karşılığını yazın.
\textbf{Ölçüt:} Aynı özellik için $\mu$--$\bar x$ veya $p$--$\hat p$ çifti
doğru tanımlanmalıdır.
\item \emph{İpucu:} Kategori, sıra, sayım ve ölçüm ayrımını kullanın.
\textbf{Örnek yanıt:} Program türü nominal, memnuniyet düzeyi sıralı, doğru
sayısı kesikli, süre sürekli değişkendir.
\item \emph{İpucu:} Kodların aritmetik anlam taşıyıp taşımadığını sorun.
\textbf{Yanıt:} Nominal kategori veya öğrenci numarası sayıyla kodlansa da
fark ve ortalamaları anlamlı değildir.
\item \emph{İpucu:} Örneklem hacmi hangi hatayı azaltır?
\textbf{Yanıt:} Büyük $n$ rastgele örnekleme değişkenliğini azaltır; kapsam,
gönüllülük veya yanıtlamama yanlılığını kendiliğinden gidermez.
\item \emph{İpucu:} Üç kümeyi iç içe düşünün.
\textbf{Ölçüt:} Hedeflenen bütün birimler, fiilen ulaşılabilen grup ve seçimde
kullanılan somut liste ayrı ayrı belirtilmelidir.
\item \emph{İpucu:} Aynı sınıftaki öğrencilerin ortak öğretmen ve çevresini düşünün.
\textbf{Yanıt:} Sınıf içi benzerlik pozitif ilişki oluşturabilir; etkin bilgi
miktarı azalır ve bağımsızlık varsayımıyla bulunan standart hata küçük kalabilir.
\item \emph{İpucu:} $\hat p=x/n$. \textbf{Yanıt:} $\hat p=145/250=0{,}58$;
parametre hedef evrendeki uygulamayı başarılı bulanların bilinmeyen oranı $p$'dir.
\item \emph{İpucu:} Birey, alt grup ve doğal küme seçimlerini karşılaştırın.
\textbf{Yanıt:} Basit rastgele seçim bireylerden; tabakalı seçim her önemli alt
gruptan; küme seçimi sınıf/okul gibi doğal gruplardan yapılır. Amaçlar sırasıyla
genel basitlik, temsil/hassasiyet ve saha maliyetini azaltmadır.
\item \emph{İpucu:} Üç eğitim düzeyini tabaka kabul edin.
\textbf{Ölçüt:} Her düzey için güncel liste, uygun tahsis kuralı ve tabaka
içinde rastgele seçim belirtilmelidir.
\item \emph{İpucu:} Ad, anlam, tür, kod ve geçerlilik bilgilerini düşünün.
\textbf{Yanıt:} Değişken adı, açıklama, ölçme düzeyi/tür, birim, kodlar,
eksik değer kodu ve geçerli aralık bilgilerinden en az beşi verilmelidir.
\item \emph{İpucu:} Yazılım 99'u gerçek değer sayarsa ne olur?
\textbf{Yanıt:} Ortalama ve dağılım bozulur; 99 eksik olarak tanımlanmalı,
analiz dosyasında uygun eksik gösterimine dönüştürülmelidir.
\end{enumerate}

\section{Betimsel istatistik}

\begin{enumerate}[leftmargin=*,itemsep=5pt]
\item \emph{İpucu:} Küçük değerlerin yanına tek bir büyük değer ekleyin.
\textbf{Örnek yanıt:} $2,3,3,4,20$ için medyan 3, ortalama 6,4'tür.
\item \emph{İpucu:} Simetri ve aykırı değerleri değerlendirin.
\textbf{Yanıt:} Yaklaşık simetrik, aykırısız dağılımda ortalama ile standart
sapma; çarpık veya aykırı değerli dağılımda medyan ile IQR tercih edilir.
\item \emph{İpucu:} Doğruluk, evrene aidiyet ve etkiyi sorgulayın.
\textbf{Ölçüt:} Veri giriş/ölçüm hatası mı, geçerli bir alt gruba mı ait,
sonuçları ne kadar değiştiriyor soruları bulunmalıdır.
\item \emph{İpucu:} Toplamı 5'e bölün ve sıralı ortanca değeri bulun.
\textbf{Yanıt:} Ortalama $5$, medyan $4$, açıklık $9-3=6$.
\item \emph{İpucu:} Ortanca dışarıda bırakılarak iki yarının ortancalarını alın.
\textbf{Yanıt:} Beş sayı özeti $(10,12,12{,}5,15,28)$, $IQR=3$; sınırlar
$7{,}5$ ve $19{,}5$ olup 28 aykırı değer adayıdır.
\item \emph{İpucu:} Ortalamayı sabit tutup merkezden uzaklıkları değiştirin.
\textbf{Örnek yanıt:} $49,50,51$ ve $40,50,60$ ortalaması 50 olan iki settir;
ikincisinin standart sapması daha büyüktür.
\item \emph{İpucu:} Her frekansı 50'ye bölün.
\textbf{Yanıt:} Yüz yüze 20 ve yüzde 40; karma 18 ve yüzde 36; çevrim içi 12
ve yüzde 24.
\item \emph{İpucu:} Değişken türü, sütunların teması ve eksen anlamını karşılaştırın.
\textbf{Yanıt:} Histogram nicel aralıkları ve bitişik sınıfları, çubuk grafiği
kategorileri ve ayrık çubukları gösterir; histogram alan/frekans dağılımını,
çubuk uzunluğu kategori büyüklüğünü temsil eder.
\item \emph{İpucu:} Uzun sağ kuyruk hangi özeti yukarı çeker?
\textbf{Yanıt:} Genellikle ortalama medyandan büyüktür; tipik gelir için
medyan ve IQR daha sağlamdır.
\item \emph{İpucu:} $200-184$ eksik gözlemi inceleyin.
\textbf{Yanıt:} Eksiklerin hangi değişken ve gruplarda oluştuğu, kodlanışı,
mekanizması ve analize etkisi kontrol edilmelidir.
\item \emph{İpucu:} Çubukların görünen uzunluklarını gerçek yüzdelerle karşılaştırın.
\textbf{Yanıt:} Kesilmiş eksen küçük farkı büyütür; sıfır tabanlı eksen,
doğrudan değer etiketi veya farkı açıkça gösteren uygun başka grafik kullanılmalıdır.
\item \emph{İpucu:} Aynı istatistikleri gözlemle ve gözlemsiz hesaplayın.
\textbf{Ölçüt:} Ortalama, medyan, standart sapma/IQR ve grafik iki durumda
karşılaştırılmalı; çıkarma kararı veri doğruluğu ve bilimsel gerekçeyle verilmelidir.
\end{enumerate}

\section{Eksik veri ve veri kalitesi}

\begin{enumerate}[leftmargin=*,itemsep=5pt]
\item \emph{İpucu:} Gözlenmemiş olmayı sayısal değer ve uygulanabilirlikle ayırın.
\textbf{Yanıt:} Boş hücre gerçek eksik olarak tanımlanır; 0 geçerli sıfırsa
korunur; 99 veri sözlüğüne göre özel kodsa eksik gösterimine çevrilir; yapısal
eksiklik ``uygulanamaz'' olarak ayrı belgelenir.
\item \emph{İpucu:} $30/200$. \textbf{Yanıt:} Eksiklik oranı yüzde 15'tir.
Genel oran, eksikliğin belirli grup veya düşük/yüksek puanlarda yoğunlaşmasını
gizleyebilir.
\item \emph{İpucu:} Eksikliğin neye bağlı olduğunu belirtin.
\textbf{Örnek yanıt:} Formların rastgele kaybı MCAR; son-test eksikliğinin
gözlenen başlangıç puanına bağlı olması MAR; iyileşmeyenlerin gözlenmeyen
sonuçları nedeniyle ayrılması MNAR örneğidir.
\item \emph{İpucu:} Miktar kadar seçiciliği düşünün.
\textbf{Yanıt:} Küçük oran kritik bir alt grupta yoğunlaşabilir; güvenlik,
yanlılık ve yöntem kararı mekanizma ile örüntüye bağlıdır.
\item \emph{İpucu:} Analize giren örneklemlerin aynı olup olmadığını sorun.
\textbf{Yanıt:} Liste bazında silme modeldeki bütün değişkenleri gözlenen tek
altküme kullanır; çift bazında silme her özet/değişken çifti için farklı
kayıtları kullanabilir.
\item \emph{İpucu:} Eksik hücrelerin hepsi merkeze ekleniyor.
\textbf{Yanıt:} Ortalama değişmeyebilir; fakat ortalamadan sıfır sapmalı yapay
gözlemler kareli sapma toplamına katkı vermeden paydayı büyütür ve varyansı
küçültür.
\item \emph{İpucu:} $T=\bar U+(1+1/m)B$.
\textbf{Yanıt:} $\bar U$ her tamamlanmış veri setindeki örnekleme belirsizliğini,
$B$ ise makul atamalar arasındaki ek belirsizliği temsil eder.
\item \emph{İpucu:} Sonucu ve eksikliği öngören gözlenen bilgileri düşünün.
\textbf{Örnek yanıt:} Başlangıç puanı, program grubu, devam oranı ve önceki
akademik başarı; zaman ve etkileşim yapısı da hedef modele göre eklenebilir.
\item \emph{İpucu:} Gözlenmeyen sonuçları sistematik kaydırın.
\textbf{Ölçüt:} Örneğin müdahale grubundaki atanan son-testlerin 3, 5 ve 8 puan
daha düşük alındığı senaryolarda tahmin, aralık ve karar karşılaştırılmalıdır.
\item \emph{İpucu:} Miktar--mekanizma--yöntem--duyarlılık sırasını kullanın.
\textbf{Ölçüt:} Değişken bazında eksik $n$/yüzde, mekanizma varsayımı, analiz
ve atama modeli, atama sayısı/birleştirme yöntemi ve en az bir duyarlılık
sonucu bulunmalıdır.
\end{enumerate}

\section{Örnekleme dağılımları}

\begin{enumerate}[leftmargin=*,itemsep=5pt]
\item \emph{İpucu:} Rastgele dalgalanma ile sistematik kaymayı ayırın.
\textbf{Yanıt:} Değişkenlik aynı yöntemle seçilen örneklemlerin farklı sonuçları;
yanlılık ise sonuçların merkezinin gerçek parametreden sistematik uzaklığıdır.
\item \emph{İpucu:} Gözlem birimi ile tekrar edilen istatistiği ayırın.
\textbf{Yanıt:} Ham dağılım bireylerin puanlarını; örnekleme dağılımı çok sayıda
örneklemden hesaplanan ortalamaları gösterir ve genellikle daha dardır.
\item \emph{İpucu:} $\Var(\bar X)=\sigma^2/n$. \textbf{Yanıt:} $n$ ile ters
orantılı azalır; örneklem hacmi $c$ katına çıkarsa varyans $1/c$ olur.
\item \emph{İpucu:} $\SE=\sigma/\sqrt n$. \textbf{Yanıt:} $18/9=2$.
\item \emph{İpucu:} Merkezdeki sistematik kaymayı düşünün.
\textbf{Yanıt:} Büyük örneklem dağılımı daraltır fakat yanlı seçim mekanizmasının
merkezini gerçek parametreye taşımaz.
\item \emph{İpucu:} $\sqrt{p(1-p)/n}$. \textbf{Yanıt:}
$\sqrt{0{,}4(0{,}6)/100}\approx0{,}049$.
\item \emph{İpucu:} İki seçimin toplamı 10 olmalıdır.
\textbf{Yanıt:} $(2,8),(4,6),(6,4),(8,2)$.
\item \emph{İpucu:} Ortak sınıf etkilerini düşünün.
\textbf{Yanıt:} Aynı öğretmen ve ortam öğrencileri benzeştirir; bağımsızlık
varsayımı zayıflar ve sıradan standart hata küçük kalabilir.
\item \emph{İpucu:} Ortalama için merkez değişmez.
\textbf{Yanıt:} $\E(\bar X)=80$, $\Var(\bar X)=16^2/64=4$,
$\SE(\bar X)=2$.
\item \emph{İpucu:} Standart hata $1/\sqrt n$ ile değişir.
\textbf{Yanıt:} Üçte bire indirmek için $n$ yaklaşık $3^2=9$ katına çıkarılmalıdır.
\item \emph{İpucu:} Oran standart hatası formülünü kullanın.
\textbf{Yanıt:} $\sqrt{0{,}6(0{,}4)/600}=0{,}02$.
\item \emph{İpucu:} Merkez ve genişliği ayrı değerlendirin.
\textbf{Yanıt:} Dağılım düşük değişkenlikli fakat yanlıdır; tahminler birbirine
yakın olsa da sistematik olarak yanlış hedef çevresindedir.
\item \emph{İpucu:} Monte Carlo hatası ile veri bilgisini ayırın.
\textbf{Yanıt:} Tekrar sayısı yalnız benzetim sonucunu kararlılaştırır;
örneklem hacmi gerçek tahminin standart hatasını ve bilgi miktarını değiştirir.
\item \emph{İpucu:} Dört sıralı çiftin ortalamasını bulun.
\textbf{Yanıt:} $(2,2),(2,6),(6,2),(6,6)$ ortalamaları $2,4,4,6$'dır;
$P(2)=1/4$, $P(4)=1/2$, $P(6)=1/4$.
\end{enumerate}

\section{Merkezi limit teoremi ve standart hata}

\begin{enumerate}[leftmargin=*,itemsep=5pt]
\item \emph{İpucu:} Ham veriyi değil hangi istatistiği düşünün.
\textbf{Yanıt:} MLT, uygun koşullarda örneklem toplamı/ortalamasının
standartlaştırılmış örnekleme dağılımını yaklaşık normale götürür.
\item \emph{İpucu:} $s/\sqrt n$. \textbf{Yanıt:} $12/6=2$.
\item \emph{İpucu:} $n$ dört katına çıkıyor.
\textbf{Yanıt:} Standart hata yarıya iner.
\item \emph{İpucu:} Bireyler ile örneklem ortalamalarını ayırın.
\textbf{Yanıt:} Standart sapma bireysel değerlerin yayılımını; standart hata
tahminin örneklemden örnekleme değişkenliğini ölçer.
\item \emph{İpucu:} MLT biçimle ilgilidir, seçim mekanizmasıyla değil.
\textbf{Yanıt:} Yanlı örneklemin ortalaması normal görünebilir fakat yanlış
merkez çevresinde yoğunlaşır; temsil sorunu sürer.
\item \emph{İpucu:} $24/\sqrt n$. \textbf{Yanıt:} $n=36$ için 4,
$n=144$ için 2; hacim dört katına çıkınca hata yarıya iner.
\item \emph{İpucu:} Hem biçim hem $1/\sqrt n$ etkisini belirtin.
\textbf{Yanıt:} $n=10$ dağılımı daha geniş ve çarpıklık taşıyabilir;
$n=100$ daha dar ve normale daha yakındır.
\item \emph{İpucu:} $np$ ve $n(1-p)$ sayılarını hesaplayın.
\textbf{Yanıt:} Beklenen başarı sayısı 4'tür; küçük olduğu için normal yaklaşım
özellikle kuyruk ve aralık hesaplarında güvenilir olmayabilir.
\item \emph{İpucu:} $\sigma$ yerine $s$ kullanıldığında ek belirsizlik vardır.
\textbf{Yanıt:} Özellikle küçük örneklemde Student $t$ dağılımına geçilir.
\item \emph{İpucu:} Birey için $\sigma$, ortalama için $\sigma/\sqrt n$ kullanın.
\textbf{Yanıt:} Bireysel $z=(72-60)/8=1{,}5$; ortalama için
$z=(64-60)/(8/4)=2$.
\item \emph{İpucu:} Önce $SE=20/8$ ve sonra standartlaştırın.
\textbf{Yanıt:} $SE=2{,}5$, $z=(105-100)/2{,}5=2$;
$P(\bar X>105)\approx0{,}0228$.
\item \emph{İpucu:} Hata üçte bire inecek.
\textbf{Yanıt:} Örneklem hacmi yaklaşık dokuz katına çıkarılmalıdır.
\item \emph{İpucu:} $25/\sqrt n\le2$ eşitsizliğini çözün.
\textbf{Yanıt:} $n\ge156{,}25$; en küçük tam sayı $157$.
\item \emph{İpucu:} Beklenen olay ve olay-dışı sayılarını hesaplayın.
\textbf{Yanıt:} $np=5$, $n(1-p)=495$; ilk sayı sınırda olduğundan normal
yaklaşım dikkatle kullanılmalı, uygun kesin/Wilson yöntemleri düşünülmelidir.
\item \emph{İpucu:} Pozitif sınıf içi ilişkiyi düşünün.
\textbf{Yanıt:} Bağımlılık etkin örneklem hacmini azaltır; $s/\sqrt n$
bağımsızlık varsayarak gerçek değişkenliği olduğundan küçük gösterebilir.
\end{enumerate}

\section{Örnekleme yöntemleri}

\begin{enumerate}[leftmargin=*,itemsep=5pt]
\item \emph{İpucu:} Hedeflenen grubun listede bulunmayan kesimini seçin.
\textbf{Örnek yanıt:} Bütün öğrenciler hedeflenirken yalnız kampüs yurdunda
kalanların listesinden seçim yapılması kapsam hatası oluşturur.
\item \emph{İpucu:} Temsil ile saha maliyetini karşılaştırın.
\textbf{Yanıt:} Tabakalama her önemli alt grubun temsilini/hassasiyeti artırır;
küme örneklemesi doğal grupları seçerek veri toplama maliyetini azaltır.
\item \emph{İpucu:} Rastgele hata ile sistematik hatayı ayırın.
\textbf{Yanıt:} Büyük $n$ rastgele dalgalanmayı azaltır; yanlı çerçeve,
yanıtlamama veya gönüllülüğün oluşturduğu sistematik farkı gidermez.
\item \emph{İpucu:} $\pi=n/N$, $w=1/\pi$.
\textbf{Yanıt:} $\pi=250/5000=0{,}05$, temel ağırlık $w=20$.
\item \emph{İpucu:} $k=N/n$ ve başlangıca ardışık $k$ değerleri ekleyin.
\textbf{Yanıt:} $k=20$; ilk beş sıra $12,32,52,72,92$.
\item \emph{İpucu:} Toplam evren 1000'dir.
\textbf{Yanıt:} Orantılı hacimler sırasıyla $210,60,30$.
\item \emph{İpucu:} Seçilme olasılığının bilinip bilinmediğini sorun.
\textbf{Yanıt:} Tabakalı örnekleme her tabakada olasılıklı seçim yapar; kota
örneklemesinde sayılar hedeflense de birimlerin seçilme olasılıkları bilinmez.
\item \emph{İpucu:} $\DEFF\approx1+(m-1)\rho$.
\textbf{Yanıt:} $1+29(0{,}05)=2{,}45$.
\item \emph{İpucu:} $n_{\rm etkin}=n/\DEFF$.
\textbf{Yanıt:} $600/2{,}5=240$.
\item \emph{İpucu:} Yanıt oranı ve tersi.
\textbf{Yanıt:} $520/800=0{,}65$; düzeltme çarpanı $1/0{,}65\approx1{,}54$.
\item \emph{İpucu:} Ağırlığın hangi değişkenlere göre kurulduğunu düşünün.
\textbf{Yanıt:} Yanıt olasılığını açıklayan değişkenler modele girmemişse veya
yanıtlamama sonuçla doğrudan ilişkiliyse grup içi fark sürer.
\item \emph{İpucu:} Her aşamada seçim birimini küçültün.
\textbf{Ölçüt:} İller tabakalanıp seçilir; illerde okullar, okullarda sınıflar,
sınıflarda öğrenciler olasılıklı seçilir; her aşama olasılığı kaydedilir.
\end{enumerate}

\section{Nokta tahmini}

\begin{enumerate}[leftmargin=*,itemsep=5pt]
\item \emph{İpucu:} Kural ile gözlenen sayıyı ayırın.
\textbf{Yanıt:} Tahmin edici örneklemden örnekleme değişen rassal kuraldır;
tahmin, belirli örneklemde aldığı sayısal değerdir.
\item \emph{İpucu:} Uzun dönem merkezi kadar saçılmayı da düşünün.
\textbf{Yanıt:} Yansız olsa da örneklemler arasında çok değişir; tek çalışmada
gerçek parametreden uzak bir sonuç verme olasılığı yüksektir.
\item \emph{İpucu:} Aynı tahmin farklı örneklem hacimlerinde ne kadar güvenilir?
\textbf{Yanıt:} Nokta tahmini belirsizliği göstermez; standart hata veya güven
aralığı ve örneklem hacmiyle verilmelidir.
\item \emph{İpucu:} Oran standart hatası.
\textbf{Yanıt:} $\sqrt{0{,}4(0{,}6)/100}\approx0{,}049$.
\item \emph{İpucu:} Hassasiyet ile doğruluğu ayırın.
\textbf{Yanıt:} Yanlı bir seçim mekanizması çok dar fakat yanlış merkezli bir
örnekleme dağılımı üretebilir.
\item \emph{İpucu:} $X$ olumlu yanıt sayısıdır.
\textbf{Yanıt:} Parametre evren oranı $p$; tahmin edici $\hat p=X/120$;
gerçekleşmiş tahmin $78/120=0{,}65$.
\item \emph{İpucu:} $MSE=\Var+\text{yanlılık}^2$.
\textbf{Yanıt:} $5+(-1)^2=6$.
\item \emph{İpucu:} İki MSE'yi ayrı hesaplayın.
\textbf{Yanıt:} $MSE_A=12$, $MSE_B=4+2^2=8$; kare hata ölçütünde B daha iyidir.
\item \emph{İpucu:} Ortalama 6'dır.
\textbf{Yanıt:} Kareli sapmalar toplamı 20; $s^2=20/(4-1)=20/3\approx6{,}67$.
\item \emph{İpucu:} Tutarlılık asimptotik bir özelliktir.
\textbf{Yanıt:} Küçük $n$'de yüksek varyans, yanlılık veya aykırı değerlere
duyarlılık sürebilir; yakınsama yalnız $n$ büyüdükçe garanti edilir.
\item \emph{İpucu:} Bütün değerleri kullanan ölçü ile sıralama ölçüsünü ayırın.
\textbf{Yanıt:} Ortalama uç değerlerden güçlü etkilenir; medyan sıralamaya
dayandığı için daha sağlamdır, ancak simetrik normal veride daha az etkin olabilir.
\item \emph{İpucu:} Grup toplamlarını birleştirin.
\textbf{Yanıt:} $[30(60)+70(75)]/100=70{,}5$.
\item \emph{İpucu:} Belirsizlik hangi ondalık basamağı destekliyor?
\textbf{Yanıt:} $SE=0{,}08$ iken dört ondalık, verinin desteklemediği kesinlik
izlenimi verir; tahmin benzer basamakta yuvarlanmalıdır.
\item \emph{İpucu:} Aynı veri üretim sürecini çok kez taklit edin.
\textbf{Ölçüt:} Parametre ve dağılımı belirleme, çok örnek üretme, her örnekte
tahmin hesaplama, ortalama eksi parametreyle yanlılık ve standart sapmayla SE
hesaplama, farklı $n$ ve koşulları karşılaştırma adımları bulunmalıdır.
\end{enumerate}

\section{Güven aralıkları}

\begin{enumerate}[leftmargin=*,itemsep=5pt]
\item \emph{İpucu:} Kritik değer nasıl değişir?
\textbf{Yanıt:} Yüzde 99 kritik değeri daha büyük olduğundan aralık genişler.
\item \emph{İpucu:} Hata payı $1/\sqrt n$ ile değişir.
\textbf{Yanıt:} Örneklem hacmi dört katına çıkarsa hata payı yaklaşık yarıya iner.
\item \emph{İpucu:} Referans değer ve olası etki büyüklüklerini ayrı okuyun.
\textbf{Yanıt:} Aralığın sıfır/referans değeri içerip içermemesi istatistiksel
kanıtı; genişliği ve uçları makul etkilerin uygulamalı büyüklüğünü gösterir.
\item \emph{İpucu:} $SE=15/5$ ve $t=2{,}064$.
\textbf{Yanıt:} Hata payı $6{,}192$; yaklaşık yüzde 95 GA
$[65{,}81;78{,}19]$.
\item \emph{İpucu:} Aralık yalnız örnekleme hatasını kapsar.
\textbf{Örnek yanıt:} Büyük fakat gönüllü/yanlı örneklem çok dar, buna rağmen
hedef evren için sistematik olarak yanıltıcı bir aralık üretebilir.
\item \emph{İpucu:} $SE=14/7=2$.
\textbf{Yanıt:} $60\pm2{,}01(2)=[55{,}98;64{,}02]$.
\item \emph{İpucu:} $\hat p=310/500$.
\textbf{Yanıt:} $\hat p=0{,}62$, $SE\approx0{,}0217$; yaklaşık yüzde 95 GA
$[0{,}577;0{,}663]$.
\item \emph{İpucu:} Tahmin $\pm zSE$.
\textbf{Yanıt:} Yüzde 90 GA $[19{,}07;28{,}94]$; yüzde 95 GA
$[18{,}12;29{,}88]$; ikincisi daha geniştir.
\item \emph{İpucu:} Hata payı $2(1{,}5)$.
\textbf{Yanıt:} GA $[1;7]$; sıfırı dışladığı için aynı çift yönlü testte
yüzde 5 düzeyinde fark yönünde kanıt vardır.
\item \emph{İpucu:} $n=(1{,}96\sigma/m)^2$.
\textbf{Yanıt:} $n\approx245{,}86$; yukarı yuvarlanarak en az 246.
\item \emph{İpucu:} Geçerli sayı beklenen yanıt oranına bölünür.
\textbf{Yanıt:} $240/0{,}75=320$ kişi davet edilmelidir.
\item \emph{İpucu:} Kapsama hedefinin maliyetini düşünün.
\textbf{Yanıt:} Daha yüksek kapsama için daha büyük kritik değer gerekir;
aynı veriyle bu, daha geniş ve daha az hassas aralık üretir.
\item \emph{İpucu:} Sıfırı ve olası iki yönü inceleyin.
\textbf{Yanıt:} Aralık sıfırı içerdiğinden yüzde 5 çift yönlü test anlamlı
değildir; eksi 2 ile artı 12 arasında iki yön ve büyük pozitif etkiler hâlâ
makuldür, dolayısıyla eşitlik kanıtlanmaz.
\item \emph{İpucu:} İki tahminin kovaryansı farkın standart hatasına girer.
\textbf{Yanıt:} Ayrı aralıkların örtüşmesi doğrudan farkın örnekleme
dağılımını vermez; karar $\mu_1-\mu_2$ için kurulan aralıktan yapılmalıdır.
\end{enumerate}

\ifimoders\else
\section{Bootstrap ve rastgeleleştirme}

\begin{enumerate}[leftmargin=*,itemsep=5pt]
\item \emph{İpucu:} Ampirik dağılımı olası evren gibi düşünün.
\textbf{Yanıt:} Yerine koyma, örneklemdeki her gözlemin her çekilişte yeniden
seçilebilmesini sağlayarak örneklem büyüklüğünü korur ve ampirik dağılımdan
bağımsız çekilişleri taklit eder.
\item \emph{İpucu:} Biri veriden, diğeri kuramsal evrenden düşünülür.
\textbf{Yanıt:} Bootstrap dağılımı gözlenen örneklemin ampirik dağılımından
üretilen yaklaşımdır; kuramsal örnekleme dağılımı gerçek veri üretim modelinin
tüm olası örneklemlerine dayanır.
\item \emph{İpucu:} Hesaplama kararlılığı ile veri bilgisini ayırın.
\textbf{Yanıt:} Büyük $B$ Monte Carlo dalgalanmasını azaltır; büyük $n$
tahminin gerçek örnekleme belirsizliğini ve bilgi miktarını değiştirir.
\item \emph{İpucu:} İki konuma 3 veya 7 gelebilir.
\textbf{Yanıt:} $(3,3),(3,7),(7,3),(7,7)$.
\item \emph{İpucu:} Bootstrap tahminlerinin yayılımını düşünün.
\textbf{Yanıt:} Tekrarlardaki $\hat\theta^*$ değerlerinin standart sapması,
$\hat\theta$ tahmin edicisinin standart hatasını yaklaşıklar.
\item \emph{İpucu:} Toplam kuyruk olasılığı yüzde 10'dur.
\textbf{Yanıt:} Yüzde 5 ve yüzde 95 yüzdelikleri kullanılır.
\item \emph{İpucu:} Yeniden örneklenen veri hangi seçim sürecinden geldi?
\textbf{Yanıt:} Bootstrap gözlenen örneklemin temsil sorununu miras alır;
eksik grupları veya sistematik seçim kaymasını oluşturup düzeltemez.
\item \emph{İpucu:} Gözlenen sonuçları sabit tutup grup üyeliklerini değiştirin.
\textbf{Yanıt:} Sonuçlar birleştirilir, özgün grup büyüklükleri korunarak
etiketler rastgele dağıtılır, fark hesaplanır ve işlem çok kez yinelenir.
\item \emph{İpucu:} Artı bir düzeltmesini kullanın.
\textbf{Yanıt:} $(124+1)/(4999+1)=125/5000=0{,}025$.
\item \emph{İpucu:} Mutlak değer ile yalnız pozitif kuyruğu karşılaştırın.
\textbf{Yanıt:} İki yönlü testte $|T^*|\ge|T_{\text{göz}}|$; üst yönlü testte
$T^*\ge T_{\text{göz}}$ sonuçları uç kabul edilir.
\item \emph{İpucu:} Eşleşme hangi birim içinde korunmalı?
\textbf{Yanıt:} Bağımsız etiket karıştırma çift bağını yok eder; farkların
işareti çift içinde değiştirilmelidir veya tasarıma uygun eşleştirilmiş yöntem kullanılmalıdır.
\item \emph{İpucu:} Bağımsız örnekleme birimi okul mu satır mı?
\textbf{Yanıt:} Aynı okuldaki satırlar bağımlıdır; satır bootstrap'i kümeyi
parçalar ve belirsizliği küçültebilir. Okullar küme olarak yeniden örneklenmelidir.
\item \emph{İpucu:} Yüzdelik aralığın simetri ve yanlılık düzeltmesini düşünün.
\textbf{Yanıt:} Küçük örneklem ve güçlü çarpıklıkta bootstrap dağılımı
parametre dağılımına kötü yaklaşabilir; basit yüzdelik aralık kapsama hatası verebilir.
\item \emph{İpucu:} Farkı yüzde puanla, testi sıfır hipotezi altında yorumlayın.
\textbf{Yanıt:} Program grubunun iyileşme oranı yaklaşık 40,5 yüzde puan daha
yüksektir; $p=0{,}0267$ etiketlenebilirlik altında bu kadar uç farkın seyrek
olduğunu gösterir. Rastgele atama/tasarım bilgisi olmadan nedensellik kurulmaz.
\item \emph{İpucu:} Belirsizlik tahmini ile hipotez sınamasını ayırın.
\textbf{Yanıt:} Bootstrap gözlenen tahminin örnekleme belirsizliğini ve aralığını;
rastgeleleştirme sıfır hipotezi altında gözlenen kadar uç sonucun olasılığını yanıtlar.
\item \emph{İpucu:} Yeniden üretim için tüm hesaplama kararlarını düşünün.
\textbf{Yanıt:} Veri kaynağı/sürümü, istatistik, $n$, tekrar sayısı $B$, tohum,
yeniden örnekleme birimi, uçluk tanımı, yazılım/paket sürümleri ve koddan en az
beşi verilmelidir.
\end{enumerate}
\fi

\section{Hipotez testlerinin mantığı}

\begin{enumerate}[leftmargin=*,itemsep=5pt]
\item \emph{İpucu:} İddiayı evren parametresi üzerinden yazın.
\textbf{Örnek yanıt:} $H_0:\mu_A-\mu_B=0$, $H_1:\mu_A-\mu_B\ne0$.
\item \emph{İpucu:} Hesap koşulu $H_0$'ın doğru kabul edilmesidir.
\textbf{Yanıt:} \pval-değeri, $H_0$ altında gözlenen kadar veya daha uç verinin
olasılığıdır; $H_0$'ın doğru olma olasılığı, etki büyüklüğü veya yeniden üretim
olasılığı değildir.
\item \emph{İpucu:} Kanıt yokluğu ile yokluk kanıtını ayırın.
\textbf{Yanıt:} ``$H_0$'ı reddetmek için yeterli kanıt bulunamadı'' denmelidir.
\item \emph{İpucu:} Aynı $p$ iki eşikle karşılaştırılır.
\textbf{Yanıt:} $p\le0{,}01$ ise ikisinde; $0{,}01<p\le0{,}05$ ise yalnız
$\alpha=0{,}05$'te reddedilir; $p>0{,}05$ ise ikisinde reddedilemez.
\item \emph{İpucu:} Örneklem hacminin \pval-değerine etkisini düşünün.
\textbf{Yanıt:} Çok küçük etkiler büyük $n$ ile anlamlı olabilir; uygulamalı
önem etki tahmini, güven aralığı ve bağlamsal eşikle değerlendirilmelidir.
\item \emph{İpucu:} Alt sınır $0{,}09$ puandır.
\textbf{Yanıt:} Pozitif fark için kanıt vardır fakat makul en küçük fark çok
küçüktür; programın eğitimsel olarak önemli etki sağladığı garanti edilmez.
\item \emph{İpucu:} Eşik karar kuralıdır, kanıtta sıçrama yaratmaz.
\textbf{Yanıt:} İki değer neredeyse aynı veri uyumsuzluğunu gösterir; 0,05'in
iki yanında olmaları bilimsel kanıtı keskin biçimde farklılaştırmaz.
\item \emph{İpucu:} Yön veri görülmeden önce gerekçelendirilmelidir.
\textbf{Örnek yanıt:} Programın başarıyı artıracağı önceden öngörülüyorsa
$H_0:\mu\le50$, $H_1:\mu>50$.
\item \emph{İpucu:} Azalma alt yönlü alternatiftir.
\textbf{Yanıt:} Hedef parametre yöntem evreninin ortalaması $\mu$;
$H_0:\mu\ge\mu_0$, $H_1:\mu<\mu_0$.
\item \emph{İpucu:} Farkı standart hataya bölün.
\textbf{Yanıt:} $Z=(62-58)/2=2$; üst yönlü $p\approx0{,}0228$, çift yönlü
$p\approx0{,}0456$.
\item \emph{İpucu:} $|Z|$ değerini 1,96 ile karşılaştırın.
\textbf{Yanıt:} $2{,}20>1{,}96$; yüzde 5 çift yönlü testte $H_0$ reddedilir.
\item \emph{İpucu:} 0,032'yi iki eşikle karşılaştırın.
\textbf{Yanıt:} $\alpha=0{,}05$'te reddedilir; $\alpha=0{,}01$'de reddedilemez.
Kanıt değişmez, yalnız karar eşiği değişir.
\item \emph{İpucu:} Aralık sıfır ve 4 puan sınırına göre nerede?
\textbf{Yanıt:} Pozitif etki için istatistiksel kanıt vardır; fakat aralığın
tamamı 4 puanın altında olduğundan etki belirlenen uygulamalı önem düzeyine ulaşmaz.
\item \emph{İpucu:} Aralığın sıfır dışındaki değerlerini de okuyun.
\textbf{Yanıt:} Sıfır uyumludur fakat eksi 1'den artı 13'e kadar, önemli
pozitif etkiler dâhil geniş bir küme de uyumludur; çalışma hassas değildir.
\item \emph{İpucu:} Testte standart hatayı $p_0$ ile hesaplayın.
\textbf{Yanıt:} $SE_0=\sqrt{0{,}4(0{,}6)/400}\approx0{,}0245$;
$Z=(0{,}45-0{,}40)/0{,}0245\approx2{,}04$.
\item \emph{İpucu:} $SE$ yaklaşık $1/\sqrt n$'dir.
\textbf{Yanıt:} Büyük $n$ aynı etki için daha küçük standart hata, daha büyük
mutlak test istatistiği ve genellikle daha küçük \pval-değeri üretir.
\end{enumerate}

\section{Hata türleri, güç ve tek örneklem testi}

\begin{enumerate}[leftmargin=*,itemsep=5pt]
\item \emph{İpucu:} Yanlış alarm ve kaçırmayı düşünün.
\textbf{Örnek yanıt:} Etkisiz programı etkili sanmak Tip I; etkili programı
belirleyememek Tip II hatadır.
\item \emph{İpucu:} Sinyal, gürültü ve karar eşiğini düşünün.
\textbf{Yanıt:} Etki büyüklüğü, örneklem hacmi, değişkenlik, $\alpha$ ve test
yönü gücü etkiler; bunlardan üçü yeterlidir.
\item \emph{İpucu:} Tek evren ortalamasını sabit bir değerle karşılaştırın.
\textbf{Yanıt:} Hedef, $\mu$ evren ortalamasının önceden belirlenmiş bilimsel
referans $\mu_0$'dan farklı/büyük/küçük olup olmadığını sınamaktır.
\item \emph{İpucu:} Çözümlü örnekte $s=10$, fark 5'tir.
\textbf{Yanıt:} $SE=10/10=1$, $t=5/1=5$; önceki $n=25$ durumuna göre hata
yarıya, test istatistiği iki katına çıkar.
\item \emph{İpucu:} Güç çalışma öncesi varsayımlara dayanır.
\textbf{Yanıt:} \pval-değeri gözlenen veriye ilişkindir; planlama için önemli
en küçük etki, değişkenlik, $\alpha$, yön ve hedef güç birlikte gerekir.
\item \emph{İpucu:} $\beta=1-\text{güç}$.
\textbf{Yanıt:} $n=10$ için $\beta\approx0{,}71$; $n=50$ için
$\beta\approx0{,}07$. Büyük örneklem gerçek etkiyi daha seyrek kaçırır.
\item \emph{İpucu:} $SE=8/4$.
\textbf{Yanıt:} $t=(42-38)/2=2$.
\item \emph{İpucu:} Referansın sonuçla birlikte seçilmesi karar kuralını değiştirir.
\textbf{Yanıt:} Veri sonrası seçim çoklu deneme/seçici raporlama yaratır ve
nominal Tip I hata oranını korumaz.
\item \emph{İpucu:} Risk yokken alarm ile risk varken alarm yokluğunu ayırın.
\textbf{Yanıt:} Yanlış alarm Tip I; gerçek riski gözden kaçırma Tip II hatadır.
\item \emph{İpucu:} $\beta=1-0{,}90$.
\textbf{Yanıt:} $\beta=0{,}10$; 500 gerçek etkili çalışmada yaklaşık 50 Tip II hata.
\item \emph{İpucu:} $SE=12/6$ ve $d=(\bar x-\mu_0)/s$.
\textbf{Yanıt:} $t=(74-70)/2=2$; $d=4/12\approx0{,}33$.
\item \emph{İpucu:} Gözlenen işaret önceden belirlenmiş yönle aynıysa.
\textbf{Yanıt:} Simetrik dağılımda üst yönlü \pval-değeri çift yönlü değerin
yaklaşık yarısıdır; yön veri görülmeden önce seçilmiş olmalıdır.
\item \emph{İpucu:} $[(1{,}96+0{,}842)/0{,}40]^2$.
\textbf{Yanıt:} Yaklaşık $49{,}07$; yukarı yuvarlayarak en az 50 (kesin
$t$ hesabı biraz daha büyük değer verebilir).
\item \emph{İpucu:} Sinyal/gürültü oranını koruyun.
\textbf{Yanıt:} Küçük etki test istatistiğinin beklenen uzaklığını azaltır;
aynı güce ulaşmak için standart hata, dolayısıyla daha büyük $n$ gerekir.
\item \emph{İpucu:} Test istatistiği $\sqrt n$ ile ölçeklenir.
\textbf{Yanıt:} $n=400$, $n=25$'in 16 katı olduğundan mutlak $t$ yaklaşık
4 kat büyür ve \pval-değeri çok daha küçük olur; etki büyüklüğü değişmez.
\item \emph{İpucu:} Çok büyük $n$ ve küçük $d$ seçin.
\textbf{Ölçüt:} Örneğin $d=0{,}10$, çok küçük $p$ ve dar GA verilip etkinin
istatistiksel olarak ayırt edildiği fakat önceden belirlenen uygulamalı eşiğin
altında kaldığı açıkça yazılmalıdır.
\end{enumerate}

\section{Bağımsız ve eşleştirilmiş örneklem testleri}

\begin{enumerate}[leftmargin=*,itemsep=5pt]
\item \emph{İpucu:} Farklı kişiler ile aynı kişilerin iki ölçümünü ayırın.
\textbf{Örnek yanıt:} İki ayrı sınıfın puanları bağımsız; aynı öğrencilerin
ön test--son test puanları eşleştirilmiştir.
\item \emph{İpucu:} Testin analiz birimi nedir?
\textbf{Yanıt:} Eşleştirilmiş test her çiftin $D_i$ farkına uygulanan tek
örneklem testidir; normallik ve aykırı değerler fark dağılımında incelenir.
\item \emph{İpucu:} Eşit olmayan varyans ve hacimleri düşünün.
\textbf{Yanıt:} Welch testi eşit varyans zorunluluğu olmadan doğru standart
hata ve serbestlik derecesi sağlar; eşitlikte de iyi performans gösterir.
\item \emph{İpucu:} Eşleşme sayıdan değil birim bağından gelir.
\textbf{Yanıt:} Eşit $n$ yalnız grup büyüklüğüdür; aynı birey, ikiz veya
eşleştirme ölçütüyle kurulmuş çiftler yoksa gözlemler bağımsızdır.
\item \emph{İpucu:} Farkın işaretini ters çevirin.
\textbf{Yanıt:} Tahmin ve iki uç noktanın işareti değişir:
$[L,U]$ aralığı ters yönde $[-U,-L]$ olur; bilimsel sonuç değişmez.
\item \emph{İpucu:} $\sqrt{s_1^2/n_1+s_2^2/n_2}$.
\textbf{Yanıt:} $\sqrt{144/36+196/49}=\sqrt8\approx2{,}83$.
\item \emph{İpucu:} Farkı standart hataya bölün.
\textbf{Yanıt:} $t=5/2=2{,}5$.
\item \emph{İpucu:} $SE_D=s_D/\sqrt n$.
\textbf{Yanıt:} $SE=6/4=1{,}5$, $t=3/1{,}5=2$.
\item \emph{İpucu:} $s_D^2=s_1^2+s_2^2-2rs_1s_2$.
\textbf{Yanıt:} $s_D=\sqrt{64+64-64}=8$.
\item \emph{İpucu:} Ortak bireysel düzey farkta iptal olur.
\textbf{Yanıt:} Pozitif korelasyon $s_D^2$ içindeki $-2rs_1s_2$ terimini
büyütür, fark varyansını ve standart hatayı azaltır.
\item \emph{İpucu:} Yalnız tam çiftler analize girer.
\textbf{Yanıt:} $n=34$, $sd=33$.
\item \emph{İpucu:} Welch--Satterthwaite yaklaşımını düşünün.
\textbf{Yanıt:} Serbestlik derecesi varyans ve örneklem hacimlerinden yaklaşık
hesaplanır; kesirli olması yöntemin doğal sonucudur.
\item \emph{İpucu:} Aralık sıfırı ve iki yönü içeriyor mu?
\textbf{Yanıt:} Tahmin 4 puan olsa da GA $[-1;9]$ sıfırı içerir; yüzde 5
düzeyinde fark belirlenemez ve hem küçük ters hem önemli pozitif farklar makuldür.
\item \emph{İpucu:} Gruplara kim, nasıl girdi?
\textbf{Yanıt:} Rastgele atama yoksa başlangıç farkları ve karıştırıcılar
grup farkını açıklayabilir; test yalnız ilişkiyi ayırt eder.
\item \emph{İpucu:} Paydadaki standartlaştırma farklıdır.
\textbf{Yanıt:} $d$ birleşik/grup standart sapması, $g$ küçük örneklem
düzeltmesi, $d_z$ farkların standart sapmasını kullanır; sayılar tanım olmadan
karşılaştırılamaz.
\item \emph{İpucu:} Ölçümler arası kovaryans ne olur?
\textbf{Yanıt:} Bağımsız test eşleşme bilgisini yok sayar; pozitif korelasyonda
genellikle gereğinden büyük standart hata ve daha düşük güç üretir.
\end{enumerate}

\section{Kategorik veriler ve ki-kare}

\begin{enumerate}[leftmargin=*,itemsep=5pt]
\item \emph{İpucu:} Grup büyüklükleri farklı olabilir.
\textbf{Yanıt:} Frekanslar örüntüyü grup büyüklüğüyle karıştırabilir; araştırma
sorusuna uygun satır/sütun yüzdesi koşullu oranları karşılaştırır.
\item \emph{İpucu:} $E_{ij}=R_iC_j/N$.
\textbf{Ölçüt:} Seçilen hücrenin satır ve sütun toplamları çarpılıp genel
toplama bölünmeli, gözlenen frekans kullanılmamalıdır.
\item \emph{İpucu:} Genel farkın nerede oluştuğunu arayın.
\textbf{Yanıt:} İşaret ve büyüklük, hangi hücrelerin bağımsızlık beklentisinden
fazla/az olduğunu gösterir; çoklu inceleme nedeniyle mekanik test gibi kullanılmaz.
\item \emph{İpucu:} Satır 40, sütun 30, toplam 80.
\textbf{Yanıt:} $E=40(30)/80=15$.
\item \emph{İpucu:} Test tablo üretim tasarımını değiştirmez.
\textbf{Yanıt:} Karıştırıcılar, seçim yanlılığı ve rastgele atama yokluğu
sürebilir; ilişki nedensel yön veya etkiyi kanıtlamaz.
\item \emph{İpucu:} Katkılar 1,00; 1,67; 1,00; 1,67'dir.
\textbf{Yanıt:} Toplam $\chi^2\approx5{,}33$.
\item \emph{İpucu:} Aynı kişi iki kez sayılıyor.
\textbf{Yanıt:} Gözlemler bağımsız değildir; ikili eşleştirilmiş sonuçta
McNemar gibi eşleşmeyi koruyan yöntem gerekir.
\item \emph{İpucu:} Kanıt gücü ile ilişki büyüklüğünü ayırın.
\textbf{Yanıt:} \pval-değeri örneklem hacminden etkilenir; Cramér $V$ ilişkinin
standartlaştırılmış büyüklüğünü verir.
\item \emph{İpucu:} Tek değişken--beklenen dağılım, aynı örneklemde iki değişken,
farklı gruplarda tek dağılım ayrımını kullanın. \textbf{Ölçüt:} Bu üç yapıya
uygun birer soru sırasıyla uyum iyiliği, bağımsızlık ve homojenlikle eşleşmelidir.
\item \emph{İpucu:} $50(80)/200$.
\textbf{Yanıt:} Beklenen frekans 20.
\item \emph{İpucu:} $(r-1)(c-1)$.
\textbf{Yanıt:} $(3-1)(4-1)=6$.
\item \emph{İpucu:} Küçük boyut $\min(1,2)=1$.
\textbf{Yanıt:} $V=\sqrt{12/(300\cdot1)}=0{,}20$.
\item \emph{İpucu:} 5'in altındaki hücre oranını bulun.
\textbf{Yanıt:} Dört hücrenin biri, yani yüzde 25'i 5'in altındadır ve en
küçük değer 3'tür; Pearson yaklaşımı kuşkuludur, Fisher/kesin yöntem düşünülmelidir.
\item \emph{İpucu:} Yalnız uyumsuz çiftler değişimi belirler.
\textbf{Yanıt:} Aynı kişilerin iki yanıtı bağımlıdır; McNemar testi başarılıdan
başarısıza ve tersine geçen uyumsuz çiftleri karşılaştırır.
\item \emph{İpucu:} İşaret verilmediğini unutmayın.
\textbf{Yanıt:} Hücre bağımsızlık beklentisinden güçlü sapar; artığın işareti
olmadan fazla mı az mı olduğu, ayrıca nedensellik veya genel etki gücü söylenemez.
\item \emph{İpucu:} Büyük $n$ küçük ilişkileri ayırt edebilir.
\textbf{Yanıt:} İlişki için istatistiksel kanıt vardır fakat $V=0{,}06$
uygulamada çok küçük olabilir; yüzdeler ve bağlamsal eşikler incelenmelidir.
\end{enumerate}

\section{Korelasyon ve basit doğrusal regresyon}

\begin{enumerate}[leftmargin=*,itemsep=5pt]
\item \emph{İpucu:} Ortak neden ve ters yönü düşünün.
\textbf{Yanıt:} Korelasyon zaman sırasını, rastgele atamayı veya karıştırıcı
denetimini sağlamaz; ters nedensellik, ortak neden ve seçim ilişki üretebilir.
\item \emph{İpucu:} $X$ ve $Y$ birimlerini cümlede kullanın.
\textbf{Ölçüt:} ``$X$ bir birim arttığında beklenen $Y$, ortalama $b_1$ birim
değişir'' biçiminde, gözlemsel veride ``ilişkilidir'' denmelidir.
\item \emph{İpucu:} Veri, model, sonuç ve tanıyı kapsayın.
\textbf{Yanıt:} Veri kaynağı/$n$, saçılım grafiği, $r$, eğim ve GA, \pval-değeri,
$R^2$, artık ve etkili gözlem kontrolleri, yazılım sürümü ve nedensellik sınırı bulunmalıdır.
\item \emph{İpucu:} ``Yüzde 96,2'si'' hangi değişkenliğe aittir?
\textbf{Yanıt:} Örneklemde son puan değişkenliğinin yüzde 96,2'si çalışma
süresiyle kurulan doğrusal model tarafından açıklanır; nedensellik/gelecek başarı garantisi değildir.
\item \emph{İpucu:} Doğrusal örüntü veri dışında sürecek mi?
\textbf{Yanıt:} Dışa taşırma gözlenmeyen bölgede aynı doğrunun geçerli olduğunu
varsayar; ilişki biçimi değişebileceğinden öngörü güvenilmez olabilir.
\item \emph{İpucu:} Artık $y-\hat y$'dir.
\textbf{Yanıt:} Pozitif artık, öğrencinin modelin altı saat için öngördüğünden
daha yüksek puan aldığını gösterir.
\item \emph{İpucu:} Bir saatlik artışı cümlede kullanın.
\textbf{Yanıt:} Model koşullarında bir saat daha uzun çalışma, ortalama son
puanda 2,78 ile 3,50 puan arasındaki artışla ilişkilidir.
\item \emph{İpucu:} Pearson $r$ yalnız doğrusal örüntüyü ölçer.
\textbf{Yanıt:} Güçlü U biçimli ilişki pozitif ve negatif doğrusal eğimleri
dengeleyip $r\approx0$ verebilir; saçılım grafiği zorunludur.
\item \emph{İpucu:} Ortalama eğri ile tek öğrencinin sapmasını ayırın.
\textbf{Yanıt:} Ortalama yanıt GA'sı belirli $X$'teki ortalama için; öngörü
aralığı yeni tek gözlem için olup artık değişkenliğini de içerdiğinden daha geniştir.
\item \emph{İpucu:} İşaret ve kareyi ayrı yorumlayın.
\textbf{Yanıt:} Orta-güçlü negatif doğrusal ilişki vardır; $R^2=0{,}36$.
\item \emph{İpucu:} $b_1=r(s_Y/s_X)$.
\textbf{Yanıt:} $b_1=0{,}50(12/4)=1{,}5$.
\item \emph{İpucu:} $b_0=\bar y-b_1\bar x$.
\textbf{Yanıt:} $b_0=40-2(10)=20$; $\hat Y=20+2X$.
\item \emph{İpucu:} Önce $\hat y$, sonra $y-\hat y$.
\textbf{Yanıt:} $\hat y=30+4(6)=54$; artık $57-54=3$.
\item \emph{İpucu:} $t=r\sqrt{(n-2)/(1-r^2)}$.
\textbf{Yanıt:} $t=0{,}35\sqrt{28/0{,}8775}\approx1{,}98$, $sd=28$.
\item \emph{İpucu:} Ortalama biçimi ile varyans biçimini ayırın.
\textbf{Yanıt:} U örüntüsü doğrusallık sorununa; huni örüntüsü sabit hata
varyansı varsayımının bozulmasına işaret eder.
\item \emph{İpucu:} Kaldıraç $X$ konumu, etki katsayı değişimiyle ilgilidir.
\textbf{Yanıt:} Uç $X$ değeri doğruya yakınsa katsayıyı az etkileyebilir;
ölçüm doğrulanmadan yüksek kaldıraç gözlemi hatalı sayılamaz.
\item \emph{İpucu:} Kanıt ile açıklanan varyansı ayırın.
\textbf{Yanıt:} Büyük örneklem eğimin sıfırdan farklılığına güçlü kanıt verir;
$R^2=0{,}08$ ise model değişkenliğin yalnız yüzde 8'ini açıklar ve öngörü zayıf olabilir.
\item \emph{İpucu:} 20 saat gözlenen aralığın çok dışındadır.
\textbf{Yanıt:} Bu ciddi dışa taşırmadır; doğrusal ilişkinin 8 saatten sonra
sürdüğüne veri yoktur ve öngörü kullanılmamalı ya da açıkça çok belirsiz sayılmalıdır.
\end{enumerate}

\ifimoders\else
\section{Çoklu doğrusal regresyon}

\begin{enumerate}[leftmargin=*,itemsep=5pt]
\item \emph{İpucu:} Ayarlanmamış ve ayarlanmış ilişkiyi ayırın.
\textbf{Yanıt:} Basit katsayı $X$ ile $Y$ arasındaki toplam doğrusal ilişkiyi;
çoklu katsayı diğer model değişkenleri sabitken koşullu ilişkiyi yanıtlar.
\item \emph{İpucu:} Aynı model değerlerindeki gözlemleri karşılaştırın.
\textbf{Yanıt:} İlgili $X_j$ dışındaki açıklayıcılar aynı kabul edilerek
$X_j$'de bir birimlik farkla ilişkili ortalama $Y$ farkı yorumlanır.
\item \emph{İpucu:} Değerleri denklemde yerine koyun.
\textbf{Yanıt:} $30+4(3)-2(5)=32$.
\item \emph{İpucu:} Yeni değişken hem $X$ hem $Y$ ile ilişkili olabilir.
\textbf{Yanıt:} Önceki katsayı karıştırılmış ortak değişkenliği içeriyorsa yeni
değişken bu kısmı ayırır; katsayının büyüklüğü hatta işareti değişebilir.
\item \emph{İpucu:} $k$ kategori için $k-1$ gösterge.
\textbf{Yanıt:} Üç kategori için iki gösterge gerekir; katsayılar seçilen
referans okul türüne göre ortalama farkları verir.
\item \emph{İpucu:} Referansı ve diğer değişkenleri belirtin.
\textbf{Ölçüt:} İlgili kategorinin beklenen $Y$ değeri, diğer açıklayıcılar
sabitken referans kategoriden ortalama 6 birim düşüktür denmelidir.
\item \emph{İpucu:} Etkileşim terimini iki $Z$ değerinde yerine koyun.
\textbf{Yanıt:} $Z=0$ için $X$ eğimi $\beta_1$; $Z=1$ için
$\beta_1+\beta_3$.
\item \emph{İpucu:} Ana katsayının diğer değişken sıfırken tanımlandığını hatırlayın.
\textbf{Yanıt:} Etkileşim varken $X$ etkisi $Z$'ye göre değişir; $\beta_1$
yalnız $Z=0$ koşulundaki eğimdir, tüm gözlemler için tek genel etki değildir.
\item \emph{İpucu:} Merkezlenmiş sıfır hangi gerçek değerdir?
\textbf{Yanıt:} Açıklayıcı ortalamasından çıkarılırsa sabit terim diğer
değişkenler ortalama düzeydeyken beklenen $Y$ olur; eğim değişmez.
\item \emph{İpucu:} $1-(1-R^2)(n-1)/(n-K-1)$.
\textbf{Yanıt:} $1-0{,}5(79/74)\approx0{,}466$.
\item \emph{İpucu:} Genel ortak katkı ile tekil standart hataları ayırın.
\textbf{Yanıt:} Açıklayıcılar birlikte modele katkı verir; korelasyonları
nedeniyle tek katsayıların belirsizliği yüksek olabilir. Anlamsız tek katsayı
``etkisiz değişken'' sonucunu tek başına vermez.
\item \emph{İpucu:} $1/(1-R_j^2)$.
\textbf{Yanıt:} $VIF=1/(1-0{,}80)=5$.
\item \emph{İpucu:} Aynı bilgiyi taşıyan açıklayıcıları düşünün.
\textbf{Yanıt:} Çoklu bağlantı katsayı standart hatalarını büyütür, işaret ve
büyüklükleri örnekleme duyarlı hâle getirir; öngörü mutlaka aynı ölçüde bozulmayabilir.
\item \emph{İpucu:} Basit regresyondaki tanıları kullanın.
\textbf{Yanıt:} Eğri, ortalama yapısında eksik doğrusal olmayan terime;
huni, sabit varyans varsayımının bozulmasına işaret eder.
\item \emph{İpucu:} Ayarlama yalnız ölçülen ve doğru modellenen değişkenler içindir.
\textbf{Yanıt:} Ölçülmemiş karıştırıcı, seçim yanlılığı, ters nedensellik ve
model hatası sürebilir; çoklu regresyon rastgele atama yaratmaz.
\item \emph{İpucu:} Seçim belirsizliği ve aşırı uyumu düşünün.
\textbf{Yanıt:} Tekrarlı \pval-elemesi Tip I hata/seçim yanlılığı ve kararsız
modeller üretir; kuramı, öngörü performansını, çoklu bağlantıyı ve uygulamalı
önemi göz ardı eder.
\end{enumerate}

\section{Tek yönlü varyans analizi}

\begin{enumerate}[leftmargin=*,itemsep=5pt]
\item \emph{İpucu:} Bir nicel yanıt ve üç veya daha fazla bağımsız grup kurun.
\textbf{Ölçüt:} Örneğin üç öğretim yöntemindeki ortalama başarı puanlarının
karşılaştırıldığı tutarlı bir araştırma sorusu verilmelidir.
\item \emph{İpucu:} Bütün evren ortalamalarını eşitleyin.
\textbf{Yanıt:} $H_0:\mu_1=\mu_2=\mu_3$; $H_1$: en az bir ortalama farklıdır.
\item \emph{İpucu:} Ayrı testlerin ortak yanlış alarm riskini düşünün.
\textbf{Yanıt:} Düzeltmesiz ikili testler aile düzeyindeki Tip I hata riskini
artırır; genel test ve düzeltilmiş karşılaştırmalar kullanılmalıdır.
\item \emph{İpucu:} $k-1$, $N-k$ ve $N-1$.
\textbf{Yanıt:} Serbestlik dereceleri sırasıyla 4, 115 ve 119'dur.
\item \emph{İpucu:} Kareler toplamlarını serbestlik derecelerine bölün.
\textbf{Yanıt:} $MS_{\mathrm{gruplar}}=60$,
$MS_{\mathrm{hata}}=15$ ve $F=4$.
\item \emph{İpucu:} Hesaplanan değeri kritik değerle karşılaştırın.
\textbf{Yanıt:} $4>2{,}866$ olduğundan $H_0$ reddedilir; en az bir ortalamanın
farklı olduğuna ilişkin yüzde 5 düzeyinde kanıt vardır.
\item \emph{İpucu:} $SS_{\mathrm{gruplar}}/SS_{\mathrm{toplam}}$.
\textbf{Yanıt:} $\eta^2=180/720=0{,}25$.
\item \emph{İpucu:} Omega-kare hata değişkenliğiyle düzeltme yapar.
\textbf{Yanıt:} Omega-kare eta-karenin küçük örneklemdeki yukarı yönlü
yanlılığını azaltmaya çalıştığı için genellikle daha küçüktür.
\item \emph{İpucu:} Genel test farkın yerini göstermez.
\textbf{Yanıt:} Tukey HSD ikili farkları aile düzeyindeki hata oranını
denetleyen eşzamanlı aralıklar ve düzeltilmiş \pval-değerleriyle inceler.
\item \emph{İpucu:} Aralık sıfırı içeriyor mu?
\textbf{Yanıt:} $[-1{,}2;4{,}8]$ sıfırı içerdiğinden bu çift için fark
belirlenemez; ortalamaların eşitliği kanıtlanmış değildir.
\item \emph{İpucu:} Hacimleri ve standart sapmaları karşılaştırın.
\textbf{Yanıt:} Eşit hacimler ile birbirine yakın 5, 6 ve 5 standart
sapmaları klasik ANOVA açısından belirgin varyans sorunu göstermez.
\item \emph{İpucu:} En küçük grubun yayılımına dikkat edin.
\textbf{Yanıt:} Hacim ve varyans dengesizliği, özellikle küçük grubun büyük
varyansı nedeniyle Welch ANOVA'yı uygun kılar.
\item \emph{İpucu:} Aynı öğrenci üç kez ölçülmüştür.
\textbf{Yanıt:} Bağımsızlık sağlanmaz; tekrarlı ölçümler ANOVA'sı veya
bağımlılığı modelleyen başka bir yöntem gerekir.
\item \emph{İpucu:} Sıra testi yalnız medyanı sınamayabilir.
\textbf{Yanıt:} Dağılım biçimleri ve yayılımları farklıysa Kruskal--Wallis
sonucu genel dağılım farkını yansıtabilir; doğrudan medyan farkı denemez.
\item \emph{İpucu:} $0{,}309>0{,}05$.
\textbf{Yanıt:} $H_0$ reddedilemez; veri en az bir ortalama farkını
belirlemek için yeterli kanıt sunmamıştır, eşitlik kanıtlanmamıştır.
\item \emph{İpucu:} Gruplara atanma sürecini değerlendirin.
\textbf{Yanıt:} Rastgele atama yoksa seçim ve karıştırıcı değişkenler grup
farkını açıklayabilir; anlamlı ANOVA tek başına nedensel etki göstermez.
\end{enumerate}
\fi

\section{Genel değerlendirme ve yöntem seçimi}

\subsection*{Genel tekrar soruları}

\begin{enumerate}[leftmargin=*,itemsep=5pt]
\item \emph{İpucu:} Yanlılık dağılımın merkezidir.
\textbf{Yanıt:} Büyük örneklem rastgele yayılımı daraltır; sistematik olarak
yanlış çerçeve veya katılım mekanizmasının merkezini düzeltmez.
\item \emph{İpucu:} MLT'nin ham veriye değil örnekleme dağılımına etkisini yazın.
\textbf{Yanıt:} Uygun koşullarda toplam/ortalamanın standartlaştırılmış
örnekleme dağılımını yaklaşık normal yaparak normal ve $t$ çıkarımlarını destekler.
\item \emph{İpucu:} Tek aralığa olasılık yüklemeyin.
\textbf{Yanıt:} Aynı yöntem çok kez uygulandığında kurulan aralıkların yaklaşık
yüzde 95'i gerçek sabit parametreyi kapsar.
\item \emph{İpucu:} Karar tablosundaki dört hücreyi düşünün.
\textbf{Yanıt:} Tip I doğru $H_0$'ı reddetme olasılığı $\alpha$; Tip II yanlış
$H_0$'ı reddedememe olasılığı $\beta$; güç $1-\beta$'dır.
\item \emph{İpucu:} Farklı kişiler ile aynı kişileri ayırın.
\textbf{Yanıt:} Bağımsız test iki ayrı grubun ortalamasını; eşleştirilmiş test
aynı/bağlı çiftlerin fark ortalamasını sınar.
\item \emph{İpucu:} Bağımsız gözlem ve beklenen frekansları yazın.
\textbf{Yanıt:} Kategorik frekanslar, her birimin tek hücreye katkısı,
bağımsız gözlemler ve yeterli beklenen frekanslar gerekir.
\item \emph{İpucu:} Ortak neden ve ters nedensellik.
\textbf{Yanıt:} Güçlü $r$ yalnız doğrusal birlikteliği ölçer; araştırma tasarımı
karıştırıcıları ve yönü belirlemedikçe nedensellik kurulamaz.
\item \emph{İpucu:} Birim, yön ve açıklanan varyansı belirtin.
\textbf{Ölçüt:} $X$'te bir birim artışla ilişkili ortalama $b_1$ birimlik $Y$
değişimi ve modelin örneklem $Y$ değişkenliğinin $100R^2$ yüzdesini açıkladığı yazılmalıdır.
\item \emph{İpucu:} Birey ve tahmin yayılımını ayırın.
\textbf{Yanıt:} Standart sapma öğrencilerin puan farklılığını; standart hata
örneklem ortalamasının evren ortalamasını tahminindeki değişkenliği ölçer.
\item \emph{İpucu:} Büyük $n$ ile küçük etki seçin.
\textbf{Ölçüt:} Küçük \pval-değeri fakat uygulamalı eşik altında tahmin ve dar güven
aralığı içeren tutarlı bir örnek verilmelidir.
\item \emph{İpucu:} Analiz birimini yazın.
\textbf{Yanıt:} Welch testinde birimler iki bağımsız grubun bireysel puanları;
eşleştirilmiş testte her çiftin fark puanıdır.
\item \emph{İpucu:} Konum ile büyüklüğü ayırın.
\textbf{Yanıt:} Artıklar ilişkinin hangi hücrelerde ve hangi yönde yoğunlaştığını;
Cramér $V$ genel ilişkinin standartlaştırılmış büyüklüğünü gösterir.
\end{enumerate}

\subsection*{Mini deneme}

\begin{enumerate}[leftmargin=*,itemsep=5pt]
\item \emph{İpucu:} Hedef grup ve katılanları ayırın.
\textbf{Yanıt:} Evren tüm birinci sınıflar, örneklem 80 gönüllüdür; temel risk
gönüllülük/öz-seçim yanlılığıdır.
\item \emph{İpucu:} $15/\sqrt{100}$.
\textbf{Yanıt:} $SE=1{,}5$.
\item \emph{İpucu:} $SE=12/6$.
\textbf{Yanıt:} $68\pm2(2)=[64;72]$.
\item \emph{İpucu:} 0,08'i 0,05 ile karşılaştırın.
\textbf{Yanıt:} $H_0$ reddedilemez; reddetmek için yeterli kanıt bulunmadığı söylenir.
\item \emph{İpucu:} Aynı öğrenciler iki kez ölçülmüştür.
\textbf{Yanıt:} Eşleştirilmiş örneklem testi; son--ön fark puanlarının dağılımı incelenir.
\item \emph{İpucu:} Asimptotik yaklaşımın koşulu sağlanmıyor.
\textbf{Yanıt:} Pearson sonucuna güvenilmemeli; bilimsel olarak uygun kategori
birleştirme, kesin/Monte Carlo yöntem veya daha büyük örneklem değerlendirilmelidir.
\item \emph{İpucu:} İşaret ve kare.
\textbf{Yanıt:} Orta-güçlü negatif doğrusal ilişki; $R^2=0{,}36$.
\item \emph{İpucu:} Kanıt gücü ile etki büyüklüğünü ayırın.
\textbf{Yanıt:} Tahmin ve güven aralığı etkinin ne kadar küçük ve ne kadar
hassas olduğunu gösterir; \pval-değeri uygulamalı önemi ölçmez.
\item \emph{İpucu:} $SE=16/8$.
\textbf{Yanıt:} $SE=2$; yaklaşık GA $75\pm4=[71;79]$.
\item \emph{İpucu:} $SE=12/6$ ve fark 4.
\textbf{Yanıt:} $t=2$; Cohen $d=4/12\approx0{,}33$.
\item \emph{İpucu:} Fark $\pm2SE$.
\textbf{Yanıt:} $t=5/2=2{,}5$; yaklaşık yüzde 95 GA $5\pm4=[1;9]$.
\item \emph{İpucu:} $(r-1)(c-1)$ ve anlamlılığın kaynağı/büyüklüğü.
\textbf{Yanıt:} $sd=(2-1)(3-1)=2$; satır/sütun yüzdeleri ile artıklar ve
Cramér $V$ incelenmelidir.
\item \emph{İpucu:} Korelasyonun karesini alın.
\textbf{Yanıt:} $R^2=0{,}45^2=0{,}2025$; basit doğrusal model örneklemdeki
yanıt değişkenliğinin yaklaşık yüzde 20,3'ünü açıklar.
\item \emph{İpucu:} Yayılım uydurulan değerle artıyor.
\textbf{Yanıt:} Sabit hata varyansı (homoskedastisite) varsayımı sorgulanır.
\end{enumerate}

\section*{Üçüncü düzey: genişletilmiş çözüm kontrolü}

Bir çözümün yalnız sonuca değil gerekçeye de tam puan alabilmesi için şu sıra
kullanılabilir:
\begin{enumerate}
  \item hedef evren, parametre ve değişkenleri tanımlama;
  \item yöntemi veri yapısına göre seçme;
  \item bağımsızlık ve yönteme özgü koşulları belirtme;
  \item ara hesapları, tahmini ve belirsizliği gösterme;
  \item kararı etki büyüklüğü ve güven aralığıyla birlikte verme;
  \item genelleme ve nedensellik sınırını araştırma tasarımına göre yazma.
\end{enumerate}

\ifimoders
Örneğin Bölüm~\ref{ch:two-sample}, Alıştırma 13 için
\else
Örneğin Bölüm 12, Alıştırma 13 için
\fi yalnız ``anlamlı değildir'' demek eksiktir.
Tam yanıt, tahminin 4 puan olduğunu, $[-1;9]$ aralığının sıfırı içerdiğini,
verinin küçük ters farklardan önemli pozitif farklara kadar geniş bir kümeyle
uyumlu olduğunu ve sonucun eşitlik kanıtı sayılamayacağını birlikte belirtir.